\section{Related Work}
\label{sec:related-work}

\subsection{Serving Architectures and Scheduling Mechanisms}
\label{sec:related-serving}

A substantial body of systems work has improved LLM inference throughput and
latency through scheduling and memory-management mechanisms. Orca introduced
iteration-level scheduling, interleaving requests at the granularity of
individual decoding steps rather than whole requests
\citep{yu2022orca}. vLLM introduced PagedAttention, a paged
virtual-memory-style KV-cache allocator that substantially raised achievable
batch sizes and became a widely adopted serving-system reference point
\citep{kwon2023vllm}. Sarathi-Serve introduces a stall-free batching
schedule that interleaves chunked-prefill computation with ongoing decode
steps to reduce head-of-line blocking under continuous batching
\citep{agrawal2024sarathi}.
Splitwise and DistServe independently proposed disaggregating the
compute-bound prefill phase and the memory-bound decode phase onto separate
hardware pools connected by fast interconnects, each showing throughput or
cost improvements from phase-specific resource allocation
\citep{patel2024splitwise,zhong2024distserve}; Mooncake extends this
direction with a KV-cache-centric disaggregated architecture that treats a
shared, elastic KV cache -- rather than the compute engines themselves -- as
the central resource to schedule around \citep{qin2025mooncake}. Llumnix
adds dynamic, cross-instance request migration and rebalancing on top of a
serving cluster \citep{sun2024llumnix}. DynamoLLM addresses cluster-level
energy and cost efficiency by dynamically reconfiguring an inference
cluster's resource allocation under latency SLOs \citep{stojkovic2025dynamollm}.

A second cluster of work targets scheduling \emph{policy} design more
specifically. VTC (Virtual Token Counter) defines an explicit token-level
fairness objective for continuous-batching schedulers and proves a bound on
the resulting service-difference guarantee \citep{sheng2024vtc}. $S^3$
predicts each request's output length ahead of generation and uses that
prediction to pack KV-cache memory more tightly at admission time, trading
occasional misprediction-driven preemption for higher achievable batch size
\citep{jin2023s3}. \citet{efficientllmscheduling2024} instead learns a
pairwise ranking model over waiting requests to prioritize decode order --
a distinct mechanism from $S^3$'s admission-time length prediction, despite
both papers appearing at NeurIPS in consecutive years and both targeting
throughput under uncertain output length. A further group of recent systems
extends this space with additional scheduling mechanisms. FastServe
introduces iteration-level preemptive scheduling, JITServe targets
SLO-aware admission under imprecise request-length information, and PARS
learns a pairwise ranking model to prioritize requests directly
\citep{fastserve2024,jitserve2025,pars2025}. Libra performs dynamic,
flexible request partitioning for unbalanced workloads, the SLAI/RAD
family derives theoretically motivated, multi-objective (throughput- and
deadline-aware) scheduling rules, and SCORPIO targets admission control
under heterogeneous per-request SLOs
\citep{libra2026,slai2025,scorpio2025}. LMetric replaces workload-specific
tuning with a single lightweight scheduling score, Apt-Serve schedules
over an adaptive hybrid KV-cache/hidden-state representation, and ProServe
unifies multi-priority request scheduling across engine- and service-layer
components \citep{lmetric2026,aptserve2025,proserve2025}. Furthermore, very
recent 2026 systems continue to explore novel scheduling paradigms:
\citet{li2026beyond} targets P99 tail latency using UniBoost---a prediction-free
interpolation between first-come-first-served and shortest-job-first---coupled
with MemGuard cache-aware preemption; \citet{liu2026hmas} addresses workload
non-stationarity and multi-tenant quality-of-service drift through
hierarchical multi-agent scheduling; and \citet{wang2026smetric} optimizes the
fundamental load-balancing versus cache-locality trade-off for multi-step
agentic workloads using session-centric routing. Each of these
papers reports an improvement over one or more baselines on
the workload(s) chosen for its own experiments; none of them, individually
or collectively, asks whether the \emph{relative ordering} they report would
be preserved under an independently sourced workload, a different metric,
or a different operating load -- the question this benchmark is built
around (\S\ref{sec:related-position}).

\subsection{Production Workloads and Characterization}
\label{sec:related-workloads}

A parallel literature documents that LLM-serving workloads observed in
production are themselves heterogeneous. BurstGPT releases 10.31 million
traces collected from a regional deployment of Microsoft's Azure OpenAI GPT
service over 213 days, characterizing request burstiness, conversation
structure, and response-length behavior, and is now published as a full
KDD~2025 paper rather than only as a preprint \citep{burstgpt2025}. Two
official Microsoft releases document Azure's own internal LLM-inference
traffic: Splitwise analyzes production conversation and code-completion
traces collected on Azure in November 2023 \citep{patel2024splitwise}, and
DynamoLLM analyzes a longer, one-week Azure trace collected in May 2024
\citep{stojkovic2025dynamollm}. \citet{kvcache_wild2025} characterizes
KV-cache reuse patterns at a large cloud provider, finding that reuse
probability and lifespan are workload-category-dependent in ways that
synthetic workload studies had not captured. \citet{ayearinllmserving2026}
reports a one-year longitudinal characterization of production traffic
spanning workload evolution, caching behavior, and load balancing (a recent
preprint at the time of writing; \S\ref{sec:related-position}).
ServeGen instead characterizes a worldwide production inference service in
order to build a principled per-client workload \emph{generator}, motivated
explicitly by the position that no single fixed trace is representative
enough on its own \citep{servegen2026}; because it characterizes the same
underlying Alibaba Cloud platform as this project's Bailian/Qwen source, we
do not count it as an independent workload provider (\S\ref{sec:workloads}).
TraceLab characterizes a distinct workload family entirely -- coding-agent
sessions rather than single-turn or short-multi-turn chat requests -- and is
used here only as out-of-distribution context, not as one of this
project's three primary sources.

\subsection{Simulators and Benchmark Frameworks}
\label{sec:related-simulators}

Vidur is the closest existing infrastructure precedent to LSSP: a
large-scale simulation framework purpose-built for exploring LLM-inference
system configurations at a fidelity and speed unavailable from
running full experiments on real hardware \citep{vidur2024}. LLMServingSim
and its successor LLMServingSim~2.0 provide a related but architecturally
different hardware/software co-simulation infrastructure, adding
iteration-granularity simulation with algorithmic-redundancy reuse and,
in the 2.0 revision, explicit support for heterogeneous accelerators and
disaggregated serving techniques \citep{llmservingsim2024,llmservingsim2_2025}.
Two recent 2026 benchmarking frameworks expand this domain: CCL-Bench 1.0
\citep{ding2026cclbench} introduces a diagnostic, trace-based evaluation
of LLM training and serving infrastructure to record fine-grained execution
evidence of kernels and communication events, while XPerf \citep{wang2026xperf}
benchmarks serving systems under agentic workloads by employing fine-grained
trace replay and full-stack profiling to handle agentic nondeterminism.
These simulators and benchmarking frameworks are built to let researchers
explore serving-system design points cheaply or load-test systems under
agentic dependencies; none of them is built around, or reports, a systematic
benchmark of \emph{comparative scheduler ranking portability} across
independently sourced workloads, operating regions, and metrics in the sense
this paper defines. LSSP's own simulator (\S\ref{sec:methodology}) is
project-local infrastructure, not a contribution positioned against this
literature; we make no claim that its output is equivalent to production
hardware latency or throughput, only that it supports a fixed, controlled
benchmark protocol for relative comparisons (\S\ref{sec:methodology}).

\subsection{Workload Dependence and Benchmark Sensitivity}
\label{sec:related-dependence}

Prior work has already established, individually, that LLM-serving
workloads are heterogeneous (\S\ref{sec:related-workloads}) and that
individual serving mechanisms and schedulers are sensitive to batching
structure, prefill/decode balance, load level, and service-level
objectives (\S\ref{sec:related-serving}). We treat both of these as accepted
background, not as a novel contribution of this work. Closest to this
paper's motivation, \citet{papaioannou2024workload} show that a serving
system's measured performance and the relative standing of design choices
can themselves shift depending on which workload is used to evaluate them,
arguing for workload-choice awareness as a first-class methodological
concern; that result is stated for single-system evaluation. What we do not
find in the literature reviewed here, including that result, is a direct
study of whether a
\emph{comparative ordering} among a fixed panel of schedulers is itself
portable across independently sourced workloads, operating regions, and
metrics -- the distinction this paper's Introduction develops between
single-system sensitivity and comparative-ranking portability.

\subsection{Ranking Portability and Benchmark Methodology}
\label{sec:related-methodology}

The statistical methodology this benchmark uses -- Kendall's tau-b and
Spearman's rho for pairwise ranking agreement, paired comparison and
block-bootstrap confidence intervals for ranking-uncertainty
quantification, and a probability-of-correct-selection framing for
benchmark sample complexity -- draws on standard nonparametric statistics
and comparative-algorithm-ranking methodology (\S\ref{sec:methodology});
these are general statistical tools, not literature claims specific to LLM
serving, and are cited here only as the methods this benchmark's analysis
plan applies to its own, independently generated data. The closest
methodological precedent we identified is from the classical parallel-job-
scheduling literature: \citet{zakay2014resampling} partition a real
scheduler workload trace into independent per-user subtraces and recombine
them to resample varied-but-realistic workloads, evaluating how a
\emph{single} parallel job scheduler's measured performance responds to
workload resampling. LSSP's sample-complexity question
(\S\ref{sec:sample-complexity}) is the same evaluation instinct --
using resampling to ask how much workload evidence a scheduler-performance
conclusion actually rests on -- applied to a different object: the
portability of a \emph{comparative} ranking among multiple schedulers
across independently sourced workloads, rather than one scheduler's
sensitivity to resampled variants of a single trace.

\subsection{Recent \emph{Journal of Supercomputing} Context}
\label{sec:related-jsc}

Three recent \emph{Journal of Supercomputing} papers illustrate the
journal's current engagement with GPU-resident LLM-inference performance
and portability, at layers adjacent to but distinct from this benchmark's
request-scheduling focus. \citet{saia2026} address deploying LLM inference
services on Slurm-managed HPC clusters without compromising existing
cluster security policy -- a \emph{cluster job-placement} scheduling
problem, one level above the \emph{intra-server request} scheduling this
benchmark studies. \citet{carmona2026mlcompilers} and
\citet{siwinska2026autotuning} each study a different sense of
``portability'' for LLM/transformer workloads on GPU (and, in the latter
case, TPU) hardware: compiler-backend selection trade-offs and
auto-tuning-framework portability across accelerator hardware,
respectively -- portability of \emph{implementation performance across
hardware}, not portability of a \emph{comparative ranking across
workloads} at fixed hardware, which is this benchmark's object of study.
We cite these three papers to situate LSSP within the journal's live
interest in LLM-serving and GPU-inference performance, not because they
address this benchmark's specific research question.

\subsection{Position of LSSP}
\label{sec:related-position}

LSSP asks a different question than the majority of the prior work surveyed
above. Prior work generally studies a scheduler, a serving system, a
simulation environment, or a workload generator/characterization in its own
right. LSSP instead studies the portability of \emph{comparative} scheduler
rankings across independently sourced workload families, operating regimes,
and metrics, using a fixed, mechanism-diverse scheduler panel and a
preregistered protocol. To our knowledge, the LSSP benchmark is distinct in
treating the portability of comparative scheduler rankings itself as the
primary object of study across independent workload sources, operating
regimes, and metrics, while explicitly quantifying ranking uncertainty and
benchmark sample complexity. A literature search targeting exactly this
question (scheduler-ranking stability/portability, cross-workload
scheduler evaluation, benchmark sample complexity and representativeness
for LLM-serving systems, current through September 2026) did not surface
prior work that treats comparative scheduler-ranking portability itself as
its primary object across independent LLM-serving workload sources; the
closest single-scheduler-focused methodological precedent
(\citet{zakay2014resampling}, \S\ref{sec:related-methodology}) predates
the LLM-serving literature entirely and evaluates one scheduler's
sensitivity to workload resampling, not comparative ranking across
schedulers. Vidur is the closest infrastructure precedent
(\S\ref{sec:related-simulators}); LLMServingSim and its successor are
related but differently scoped simulator infrastructure; the recent
production-characterization literature
(\S\ref{sec:related-workloads}) establishes the workload heterogeneity this
benchmark's design responds to, but does not itself study cross-source
ranking portability. This wording is intentionally narrower than an
unconditional claim of priority; we do not claim to be the first to study
LLM-serving workload dependence or scheduler sensitivity, only that we did
not identify a prior benchmark treating cross-source, cross-metric,
cross-load-region scheduler-ranking portability as its primary object of
study under the primary-source review conducted for this manuscript.

A note on provider independence is warranted here because it bears
directly on how several of the sources above should be read together:
BurstGPT and this project's Azure trace both ultimately sit on Microsoft's
Azure infrastructure, collected by different teams through different
mechanisms, so we treat them as distinct, independently released workload
sources rather than as independent providers
(\S\ref{sec:sources} gives the full treatment, including citations).
